% ****************************************************************************************************
% Chapter 1: The Problem of Emergency Triage
% ****************************************************************************************************

\chapter{The Problem of Emergency Triage}
\label{ch:triage-problem}

Every minute, in emergency departments around the world, a nurse meets a patient they have never seen before and decides how long that patient can safely wait. This decision is called triage, and it is one of the most consequential judgements in medicine. A patient with a heart attack may look deceptively stable; a patient complaining loudly of pain may be in no immediate danger. The triage nurse has minutes to sort them, and the consequences of getting it wrong are measured in lives.

This report describes a machine learning system that attempts to assist that decision. Before I discuss the model, the data, or the experiments, I need to be clear about what I am trying to help with and why it is hard.

\section{What Triage Is}

The word \emph{triage} comes from the French \emph{trier}, to sort. It was formalised in the Napoleonic wars by Dominique Jean Larrey, Napoleon's chief surgeon, who decided that battlefield casualties should be treated in order of clinical urgency rather than rank or nationality \citep{nakao2017trauma}. The principle has held for two centuries: when medical resources are limited, they must be allocated by need.

In a modern emergency department (ED), triage happens at the front door. A nurse takes a brief history, measures vital signs, asks about the chief complaint, and assigns a priority level. That level determines how long the patient waits to see a physician, which treatment rooms they are allocated to, and sometimes whether laboratory tests are ordered pre-emptively. In a busy ED, the difference between the highest and lowest priority level is the difference between being seen in minutes and being seen in hours.

\begin{notebox}[A Note on Language]
Throughout this report, \emph{acuity} refers to the severity of a patient's clinical condition at presentation. Higher acuity means more urgent. I use \emph{under-triage} to mean assigning a patient a lower acuity than their condition warrants, the dangerous direction of error. \emph{Over-triage} means the opposite: assigning higher acuity than necessary. As the next section shows, these two errors have very different consequences.
\end{notebox}

\section{The Emergency Severity Index}

The dominant triage system in North American emergency departments is the Emergency Severity Index, universally abbreviated as ESI. It was developed in the late 1990s at the Harvard-affiliated Emergency Medicine Network and is now used in more than seventy percent of US EDs \citep{tanabe2004reliability}.

ESI assigns patients to one of five levels:

\begin{itemize}
 \item \textbf{ESI-1 (Resuscitation)}: patients in immediate life-threatening condition. Cardiac arrest, severe trauma, status epilepticus. They are seen immediately, by a physician, with a full resuscitation team.
 \item \textbf{ESI-2 (Emergent)}: patients in high-risk situations. Chest pain suggestive of myocardial infarction, severe respiratory distress, altered mental status. They should be seen within ten minutes.
 \item \textbf{ESI-3 (Urgent)}: patients requiring multiple resources but not in immediate danger. Moderate abdominal pain, fever with comorbidities. Typically seen within thirty to sixty minutes.
 \item \textbf{ESI-4 (Less Urgent)}: patients requiring one resource. Simple lacerations, minor infections. May wait one to two hours.
 \item \textbf{ESI-5 (Non-Urgent)}: patients requiring no resources beyond examination. Medication refills, minor cosmetic concerns. May wait several hours without harm.
\end{itemize}

The ESI algorithm is intentionally heuristic rather than algorithmic. It asks the nurse a series of branching questions (does the patient require immediate life-saving intervention? is the situation high-risk?) and relies heavily on clinical gestalt. Two trained nurses presented with the same patient can, and regularly do, reach different conclusions.

\section{The Scandinavian Alternative: RETTS}

The approach taken in Scandinavia is substantially different. The Rapid Emergency Triage and Treatment System (RETTS), originally called METTS, was developed by Bengt Widgren and Mehrsa Jourak at Sahlgrenska University Hospital in Gothenburg in the early 2000s \citep{widgren2011metts}. RETTS is now the dominant triage system across Sweden, Norway, and Denmark, used by the majority of emergency departments nationwide \citep{wireklint2021national}.

Where ESI trusts clinical gestalt, RETTS insists on structure. The system employs a two-step algorithmic approach:

\begin{enumerate}
 \item \textbf{Vital signs algorithm}: the patient's airway, breathing, circulation, disability (neurological status), and exposure (temperature) are evaluated against explicit numerical thresholds. Each component yields a color-coded priority: red (highest), orange, yellow, green, or blue (lowest).
 \item \textbf{Emergency Symptoms and Signs (ESS) algorithm}: one of 43 complaint-specific algorithms is selected based on the chief complaint. Chest pain has its own flowchart; dyspnoea has another; abdominal pain has another. Each yields an independent priority.
\end{enumerate}

The final triage level is the \emph{more urgent} of the two steps. This ``max'' function is the key design decision. A patient with normal-sounding chief complaint but abnormal vital signs will be prioritised by their vital signs. A patient with normal vital signs but a chief complaint that raises red flags will be prioritised by the complaint. RETTS structurally encodes the principle that when in doubt, escalate.

This philosophical difference matters for machine learning. ESI's reliance on gestalt means that two nurses can reasonably disagree; there is inherent noise in the label. RETTS's algorithmic structure reduces that noise but requires patients to fit its categorisation scheme. The published reliability data reflects both strengths and weaknesses, and I return to it below.

\begin{remarkbox}[Why Sweden Matters Here]
This project was undertaken for a competition sponsored by the Laitinen-Fredriksson Foundation, a Nordic organisation. The methods I develop for ESI prediction transfer directly to RETTS prediction because both are five-level ordinal scales, but the RETTS design itself, particularly its age-unadjusted vital sign thresholds, suggested an important hypothesis I explore in Chapter~\ref{ch:fairness}.
\end{remarkbox}

\section{The Danger of Under-Triage}

Triage errors are not symmetric. Assigning ESI-3 to a patient who should be ESI-2 means delaying care for someone who needed it urgently. Assigning ESI-2 to a patient who should be ESI-3 means occupying a high-priority resource that another patient may need. Both are errors, but only the first can kill.

The clinical literature is unambiguous on this point. \citet{haas2010survival} studied 11,398 severely injured patients in Ontario and found that under-triaged patients had an odds ratio of 1.24 for mortality (95\% CI 1.10--1.40). Four percent of under-triaged patients died before transfer could be arranged, representing twenty-two percent of all study deaths. \citet{rubano2016undertriage}, in a study at a Level II trauma centre, reported an even starker odds ratio of 2.18 (95\% CI 1.57--3.01) after controlling for age, injury severity, and comorbidities. A landmark New England Journal of Medicine study by \citet{mackenzie2006national} demonstrated a twenty-five percent relative mortality reduction when severely injured patients received trauma centre care rather than being under-triaged to non-trauma facilities.

The largest available study of mistriage specifically in ED settings comes from \citet{sax2023mistriage}, who analysed 5,316,832 ED encounters across 21 Kaiser Permanente emergency departments. The overall mistriage rate was 32.2\%, with 3.3\% under-triaged and 28.9\% over-triaged. More strikingly, the sensitivity of ESI for correctly identifying patients who required immediate life-stabilising interventions was only 65.9\%. One third of critically ill patients were assigned mid- or low-acuity levels.

The institutional response to this asymmetry is codified in the American College of Surgeons Committee on Trauma guidelines, which specify an under-triage target of $\leq 5\%$ while tolerating over-triage rates of 25--50\% \citep{acscot2014resources}. This 10:1 ratio is not arbitrary. It reflects the institutional position that missing a critically ill patient is an order of magnitude more dangerous than temporarily over-allocating resources to a stable one. This asymmetric standard becomes central to my own methodology in Chapter~\ref{ch:safety}.

\begin{notebox}[Why Odds Ratios Matter]
An odds ratio of 2.18 means that, holding other factors equal, under-triaged patients in that study had roughly twice the odds of dying compared to correctly-triaged patients with the same underlying condition. This is not the same as a doubled mortality rate, the baseline mortality matters, but it is a strong statistical signal that the direction of the error has clinical consequences.
\end{notebox}

\section{The Reliability Problem}

If triage decisions were perfectly reproducible, the task would merely be difficult. It is not reproducible. Two trained nurses presented with the same patient at the same moment will, surprisingly often, assign different acuity levels.

The definitive meta-analysis of ESI inter-rater reliability comes from \citet{mirhaghi2015esi}, who pooled 19 studies covering 40,579 cases across six countries. The overall pooled Cohen's kappa was $\kappa = 0.791$ (95\% CI 0.752--0.825), classified as \emph{substantial} agreement on the Landis-Koch scale. But this headline number conceals important subgroup variation. Expert-to-expert agreement reached $\kappa = 0.900$, effectively ceiling. Nurse-to-expert agreement was substantially lower at $\kappa = 0.732$. Most importantly, paper-based scenarios produced $\kappa = 0.824$, while \emph{live patient assessments} dropped to $\kappa = 0.694$.

This gap between scenario reliability and live reliability is striking. \citet{travers2009triage} documented it even more starkly in pediatric triage: weighted $\kappa = 0.77$ for scenarios but only 0.57 for live patients. The act of triaging a real patient, in real time, in a real ED, introduces noise that paper studies do not capture.

RETTS shows a similar pattern, with generally lower reliability than ESI in recent studies. \citet{widgren2011metts} reported $\kappa = 0.761$--0.903 in the original validation. But \citet{wireklint2018reliability} found $\kappa = 0.455$ among Swedish nurses using RETTS on paper scenarios, and \citet{olsson2022educational} found Fleiss $\kappa = 0.411$ among student nurses after formal training. These lower numbers have raised concern in the Swedish emergency medicine community about RETTS's scientific foundation \citep{sbu2010triage}.

\begin{notebox}[What Does Cohen's Kappa Measure?]
Cohen's kappa ($\kappa$) quantifies agreement between two raters beyond what would be expected by chance. A $\kappa$ of 1.0 means perfect agreement; 0.0 means agreement no better than random. The Landis-Koch scale labels $\kappa$ values as:

\begin{itemize}[leftmargin=2em, itemsep=0.1em]
 \item $< 0.20$: slight agreement
 \item $0.21 - 0.40$: fair
 \item $0.41 - 0.60$: moderate
 \item $0.61 - 0.80$: substantial
 \item $0.81 - 1.00$: almost perfect
\end{itemize}

For a five-level scale like ESI, achieving $\kappa > 0.7$ on live patients is genuinely difficult. My own model achieves $\kappa = 0.701$, a result I return to in Chapter~\ref{ch:results}.
\end{notebox}

\section{The Case for Machine Learning}

Given these two problems, the asymmetric risk of errors and the limited reliability of human raters, there is a clear case for machine-assisted triage. A model trained on large volumes of historical ED data could, in principle, learn patterns that nurses miss. It could apply consistent criteria across shifts, hospitals, and patient populations. It could make the asymmetry of errors explicit in its decision rule, pushing toward higher acuity when uncertain.

The case is not uncontested. Every additional layer of automation in medicine raises concerns about accountability, transparency, and the risk of systematic bias being scaled silently. The European Union has already classified emergency triage AI as \emph{high-risk} under Annex III, Section 5(d) of Regulation 2024/1689, the AI Act \citep{eu2024aiact}. Systems deployed in EU member states after August 2026 will be required to demonstrate risk management, bias monitoring, transparency, and human oversight. These requirements shape how I design and evaluate my system throughout this report.

The published literature on ML triage to date is substantial but uneven. \citet{levin2018machine} deployed a random forest system at Johns Hopkins, achieving AUC 0.73--0.92 for critical care prediction and reducing arrival-to-admission decision time by 58 minutes. \citet{raita2019emergency} compared gradient-boosted trees and deep neural networks on 135,470 NHAMCS visits. \citet{kwon2018deep} trained a deep learning system on 11.6 million Korean encounters, achieving AUROC 0.935 for in-hospital mortality.

But a striking feature of this literature, I discovered while planning my own work, is that almost none of it predicts ESI level directly. The vast majority of published ML triage models predict \emph{clinical outcomes}, mortality, ICU admission, hospitalisation, rather than the 5-level ESI classification itself. \citet{ivanov2021kate} developed KATE, one of the few direct ESI predictors, reporting 75.7\% overall accuracy but not macro F1 or kappa. The 5-class ESI prediction task at scale, on real clinical data, has been largely unexplored.

This report describes one attempt to fill that gap.

\section{What This Report Describes}

What follows is an account of a year's work building an ML-assisted ESI prediction system. The work began as an entry in a Kaggle competition, the Triagegeist competition sponsored by the Laitinen-Fredriksson Foundation. It became, through a discovery I describe in the next chapter, something substantially more ambitious.

A brief note on background may be useful here, since it shaped the project in ways that go beyond method choices. I trained as an astrophysicist and spent most of my career working with observational data: noisy, censored, contaminated by instrumental effects, and never quite matching the clean idealisations of the textbook. This is also, roughly, what emergency department data looks like. The instincts that transfer most readily from physical-science research are not about models but about measurements: every number has an uncertainty, every dataset has systematic effects worth looking for, and the honest answer is sometimes that the data cannot support the question being asked. Several of the choices described below, particularly the decision to pivot away from the competition dataset once its label leakage became clear, reflect those habits more than any specific machine-learning technique. Chapter~\ref{ch:reflections} returns to this transition in more personal terms at the end of the report.

The report is structured as a narrative. Chapter~\ref{ch:discovery} describes the competition and the synthetic data problem I identified. Chapter~\ref{ch:mimic} describes the pivot to real clinical data and the MIMIC-IV-ED dataset. Chapters~\ref{ch:features} through \ref{ch:architecture} describe the methodology: feature engineering, clinical scoring, text embeddings via fine-tuned Bio\_ClinicalBERT, and the gradient-boosted model blend. Chapter~\ref{ch:results} presents the headline results and benchmarks them against the published human reliability literature. Chapter~\ref{ch:safety} introduces the main methodological contribution: a head-to-head comparison of two strategies for encoding clinical risk tolerance. Chapter~\ref{ch:fairness} examines fairness and the unexpected finding that age, not race, is the dominant disparity axis in the model's errors. Chapter~\ref{ch:explainability} addresses explainability and EU AI Act compliance. The remaining chapters cover smaller experiments, honest limitations, and reflections on the research process.

The narrative includes experimental dead ends as well as successes. I believe this is important. A research project is not a series of confirmed hypotheses; it is a series of attempts, some of which work and many of which do not. Presenting only the successful experiments would misrepresent both the science and the experience of doing it.

I begin, in the next chapter, with the discovery that changed the shape of the project.
